\documentclass[../../../include/open-logic-section]{subfiles}

\begin{document}

\olfileid{sth}{cardinals}{hp}
\olsection{అనుబంధం: హ్యూమ్ సూత్రం}

\olref[cp]{sec}లో కాంటర్ సూత్రాన్ని వివరించాం. అది ఇదీ:
\begin{align*}
	\card{A} = \card{B} & \text{ అప్పుడూ అప్పుడే } A \approx B.
\intertext{ఇది ఇప్పుడు \emph{హ్యూమ్ సూత్రం} అని పిలిచే దానికి చాలా దగ్గరగా ఉంది. ఆ సూత్రం ఇలా చెబుతుంది:}
	\fregenum{x} {F(x)} = \fregenum{x}{G(x)} & \text{ అప్పుడూ అప్పుడే } F \sim G
\end{align*}
ఇక్కడ `$F \sim G$' అంటే $F$లు, $G$లు ఖచ్చితంగా ఒకే సంఖ్యలో ఉన్నాయని సంక్షిప్తంగా రాసినది; అంటే $F$లను $G$లతో ద్వైజెక్షన్‌లో పెట్టవచ్చు. మరింత స్పష్టంగా:
\begin{align*}
	\exists R(&\forall v\forall y(Rvy \lif (Fv \land Gy)) \land {}\\
		&\forall v(Fv \lif \lexists![y][Rvy]) \land {}\\
		&\forall y(Gy \lif \lexists![v][Rvy]))
\end{align*}
కానీ హ్యూమ్ సూత్రానికీ కాంటర్ సూత్రానికీ రకపరమైన తేడా ఉంది. కాంటర్ సూత్రంలో ``$A$'', ``$B$'' అనే చరరాశులు \emph{సమితులను} సూచించే ప్రథమ-క్రమ పదాలు. హ్యూమ్ సూత్రంలో ``$F$'', ``$G$'', ``$R$'' \emph{ప్రథమ-క్రమ పదాలు కావు}; అవి \emph{విధేయ స్థానంలో} ఉన్నాయి. (బహుశా అవి \emph{ధర్మాలను} సూచిస్తాయి.) అందువల్ల హ్యూమ్ సూత్రాన్ని ఇలా చెప్పవచ్చు: $F$ల సంఖ్య $G$ల సంఖ్యకు సమానం అప్పుడూ అప్పుడే $F$లు,~$G$ల మధ్య ద్వైజెక్షన్ ఉంటే. దీన్ని \emph{హ్యూమ్ సూత్రం} అంటారు; ఎందుకంటే హ్యూమ్ ఒకసారి ఇలా రాశారు:
\begin{quote}
  రెండు సంఖ్యలను పోల్చినప్పుడు, ఒకదానిలోని ప్రతి ఏకకానికి మరొకదానిలో ఒక ఏకకం సరిపోతే, వాటిని సమానమని ప్రకటిస్తాం.
  \citep[Pt.III Bk.1 \S1]{Hume1740}
\end{quote}
హ్యూమ్‌లోని ఈ భాగాన్ని ఉటంకించి, ఆ తరువాత (మనం హ్యూమ్ సూత్రం అని పిలిచిన) భావనను సంక్షిప్తంగా వివరించిన \citet[\S63]{Frege1884} వల్ల హ్యూమ్ సూత్రానికి ఆధునిక గణిత-తార్కిక ప్రాధాన్యం వచ్చింది.

హ్యూమ్ సూత్రానికీ మౌలిక నియమం~Vకీ నిర్మాణపరమైన పోలికను గమనించండి. \olref[story][blv]{sec}లో దాన్ని ఇలా సూత్రీకరించాం:
\[
	\fregeext{x}{F(x)} = \fregeext{x}{G(x)} \text{అప్పుడూ అప్పుడే } \lforall[x][(F(x) \liff G(x))].
\]
ఇప్పుడు కొంత వివరణ, పోలిక ఉపయోగపడతాయి.

హ్యూమ్ సూత్రం లేదా మౌలిక నియమం~V వంటి సూత్రాన్ని రెండు రకాలుగా అర్థం చేసుకోవచ్చు: \emph{స్వవర్గానవలంబితంగా} లేదా \emph{స్వవర్గావలంబితంగా} (\olref[story][predicative]{sec} గుర్తుచేసుకోండి). మౌలిక నియమం~Vను స్వవర్గావలంబితంగా చదివితే, ప్రతి~$F$కు $\fregeext{x}{F(x)}$ అనే వస్తువు అదే నియమాన్ని సూత్రీకరించడానికి వాడిన పరిమాణకరణ పరిధిలో ఉంటుంది. అలాగే హ్యూమ్ సూత్రాన్ని స్వవర్గావలంబితంగా చదివితే, ప్రతి~$F$కు $\fregenum{x}{F(x)}$ అనే వస్తువు అదే సూత్రాన్ని సూత్రీకరించడానికి వాడిన పరిమాణకరణ పరిధిలో ఉంటుంది. దీనికి భిన్నంగా \emph{స్వవర్గానవలంబిత} అర్థంలో $\fregeext{x}{F(x)}$,~$\fregenum{x}{F(x)}$ అనే వస్తువులు ఏదో \emph{వేరే} పరిధికి చెందినవి.

ఇప్పుడు మౌలిక నియమం~Vను స్వవర్గావలంబితంగా చదివితే, అమాయక అవగాహన ద్వారా అసంగతత వస్తుంది (వివరాలకు \olref[story][blv]{sec} చూడండి). అమాయక అవగాహనలాగే, దాన్ని \emph{స్వవర్గానవలంబితంగా} చదివితే సుసంగతంగా చేయవచ్చు. కానీ అప్పుడు అది మనం కోరుకున్న పనులన్నిటినీ చేయకపోవచ్చు.

అయితే హ్యూమ్ సూత్రాన్ని \emph{స్వవర్గావలంబితంగా} సుసంగతంగా చదవవచ్చు. అలా చదివితే అది చాలా శక్తివంతం.

ఉదాహరణకు, ఏ వస్తువూ సంతృప్తిపరచని ``$x \neq x$'' విధేయాన్ని తీసుకోండి. హ్యూమ్ సూత్రం అప్పుడు $\# x( x\neq
x)$ అనే వస్తువును ఇస్తుంది. దాన్ని~$0$ సంఖ్యగా తీసుకోవచ్చు. ఇప్పుడు \emph{స్వవర్గావలంబిత} అర్థంలో---\emph{అందులో మాత్రమే}---ఈ $0$ వస్తువు మన మొదటి పరిమాణకరణ పరిధిలో ఉంటుంది. కాబట్టి ``$x = 0$'' విధేయానికి హ్యూమ్ సూత్రాన్ని వర్తింపజేసి $\#x (x =
0)$ అనే వస్తువును పొందవచ్చు. దాన్ని~$1$ సంఖ్యగా తీసుకోవచ్చు. అంతేకాక $0 \neq 1$ అని హ్యూమ్ సూత్రం వల్ల వస్తుంది: తనకు తాను సమానంకాని వస్తువుల నుంచి $0$కు సమానమైన వస్తువులకు ద్వైజెక్షన్ ఉండదు (మొదటి రకానికి ఏదీ లేదు, రెండో రకానికి ఒకటి ఉంది). మళ్లీ స్వవర్గావలంబితంగా పనిచేస్తే, $1$~మన మొదటి పరిమాణకరణ పరిధిలో ఉంటుంది. కాబట్టి ``$(x = 0 \lor x = 1)$'' విధేయాన్ని హ్యూమ్ సూత్రంలో వాడి $\#x(x = 0 \lor
x = 1)$ అనే వస్తువును పొందవచ్చు. దాన్ని~$2$ సంఖ్యగా తీసుకోవచ్చు; $0\neq 2$, $1 \neq 2$ అని కూడా చూపవచ్చు; ఈ ప్రక్రియ ఇలాగే సాగుతుంది.

సంక్షిప్తంగా, హ్యూమ్ సూత్రాన్ని స్వవర్గావలంబితంగా తీసుకుంటే \emph{అనంత సంఖ్యలో వస్తువులు} ఉన్నాయని వస్తుంది. అందుకే \emph{నవ-ఫ్రేగే తార్కికవాదులు} దాన్ని అంకగణితానికి పునాదిగా తీసుకోవడానికి ప్రోత్సాహం పొందారు.

కానీ ఫ్రేగే \emph{స్వయంగా} హ్యూమ్ సూత్రాన్ని అంకగణిత పునాదిగా తీసుకోలేదు. దానికి బదులు, ఒక స్పష్టమైన నిర్వచనం నుంచి హ్యూమ్ సూత్రాన్ని నిరూపించారు: $\fregenum{x}{F(x)}$ను $F \sim \Phi$ భావన విస్తరణగా నిర్వచించారు. ఆధునిక పదాల్లో దాన్ని $\fregenum{x}{F(x)} = \Setabs{G}{F \sim G}$గా ఇవ్వడానికి ప్రయత్నించవచ్చు; అయితే అలా చేస్తే అమాయక అవగాహన సమస్యలకే తిరిగి వస్తాం.

\end{document}
